\subsubsection*{定理 1.1.}
固定特征指数 $p$、维数 $n\geq1$、整数 $r\geq1$ 及多重次数 $(d_1,\dots,d_r)$. 若
\begin{equation}\tag{1.1.0}
\sum_{i=1}^{r}d_i\geq(n+r)+1,
\end{equation}
则在特征 $p$ 的 $\mathbb P^{n+r}$ 中所有维数为 $n$、多重次数为 $(d_1,\dots,d_r)$
的完全交中,“充分一般”的成员在 1.0 的意义下是非特殊的.

\medskip
\noindent\emph{证明.}
设
\begin{equation}\tag{1.1.1}
f:X_U\longrightarrow U
\end{equation}
表示特征 $p$ 的 $\mathbb P^{n+r}$ 中多重次数为 $(d_1,\dots,d_r)$ 的完全交
连同其定义方程组成的“泛族”. 因而 $U$ 是 $\mathbb F_p$ 上 $r$ 个射影空间的乘积
中的开子集, 其维数相当大.

$\mathcal O_U$ 模
\begin{equation}\tag{1.1.2}
R^nf_*\mathcal O_{X_U}
\end{equation}
是局部自由的, 且其形成与任意基变换可交换 (XI~1.5(i)). 由 (1.1.0),
其秩 $M$ 非零 (XI~2.5). 当 $p=1$ 时, 这就对每个满足 (1.1.0) 的完全交证明了定理.

现设 $p\geq2$. 迭代的 Hasse--Witt 运算
\begin{equation}\tag{1.1.3}
\mathfrak F_p^M=\mathfrak F_{p^M}:
R^nf_*\mathcal O_{X_U}\longrightarrow R^nf_*\mathcal O_{X_U}
\end{equation}
在 $\mathcal O_U$ 上是 $p^M$ 线性的. 由于 $R^nf_*\mathcal O_{X_U}$ 局部自由,
(1.1.3) 在 $U$ 上局部由元素属于 $\mathcal O_U$ 的 $(M,M)$ 矩阵 $H$ 表示. 因此,
\begin{equation}\tag{1.1.4}
\{u\in U\mid X_u\text{ 非特殊}\}\quad\text{是下述空间的开集: }U.
\end{equation}
事实上, 这正是使 $H(u)\ne0$ 的 $u$ 的集合. 由于 $U$ 不可约,
只需找到一个 $u$ 使 $X_u$ 非特殊; 我们将在有限域上这样做.

\subsubsection*{命题 1.1.5.}
设 $X$ 是 $\mathbb F_q$ 上的完全交. 若 $X$ 没有 $\mathbb F_q$ 有理点, 则 $X$ 非特殊.

\medskip
\noindent\emph{证明.}
显然只需对完全交 $X$ 证明下述蕴涵
\begin{equation}\tag{1.1.5.1}
X\text{ 特殊}\quad\Longrightarrow\quad
\operatorname{card}\bigl(X(\mathbb F_q)\bigr)\equiv1\pmod p.
\end{equation}
由下一报告 XXII 的 3.2 中证明的迹公式,
\begin{equation*}
\operatorname{card}\bigl(X(\mathbb F_q)\bigr)
\equiv\sum_{i\geq0}(-1)^i
\operatorname{tr}\bigl(\mathfrak F_q\mid H^i(X,\mathcal O_X)\bigr)
\pmod p.
\end{equation*}
而除非 $i=0$ 或 $i=n=\dim X$, 都有 $H^i(X,\mathcal O_X)=0$,
且 $H^0(X,\mathcal O_X)\simeq\mathbb F_q$ (XI~1.5(iii)); 因而上式化为
\begin{equation*}
\operatorname{card}\bigl(X(\mathbb F_q)\bigr)
\equiv1+(-1)^n
\operatorname{tr}\bigl(\mathfrak F_q\mid H^n(X,\mathcal O_X)\bigr)
\pmod p.
\end{equation*}
若 $X$ 特殊, 则依定义 $\mathfrak F_q$ 在 $H^n(X,\mathcal O_X)$ 上的运算幂零,
故其迹为零. 这就证明了 (1.1.5.1), 从而证明了命题.

\medskip
\noindent\emph{底本注.}
照录法文层保留了印本中的 \emph{parmi toute les}、\emph{théroème}、
\emph{intersection complète vérifiante} 及 \emph{n'a pas des points};
校订法文层和译文采用底本缺陷登记表中记录的合乎语法的读法.
